\section{Rigidity of the normalized static base}
\label{hw-section}

The equality data determine the complete static base.  We use a
static Heintze--Karcher deficit on smooth bounded domains, followed
by an exhaustion and the finite-domain warped-product rigidity of
\cite[Theorem~1.1]{borghinifogagnolopinamonti2024}.

\begin{theorem}[A static Heintze--Karcher deficit]
\label{hw-public-rigidity}
Let $(M^3,q,V)$ be a smooth static system satisfying
\[
 \Hess_qV=V(\Ric_q+3q),\qquad \Delta_qV=3V,
\]
with $V>0$ in the interior and connected compact nondegenerate
horizon boundary $S=\{V=0\}$.  Suppose a bounded smooth domain
$\Omega$ has boundary $S\sqcup\Sigma$, where $\Sigma$ has strictly
positive mean curvature $H$ toward the exterior of $\Omega$.
Set $A=\Area_q(S)$ and $\kappa=|\nabla V|_S$.  Then
\[
 D=3A+2\pi\chi(S)>0,\qquad c=\frac{2A}{3D}>0.
\]
The unique smooth solution of
\begin{equation}\label{hw-deficit-dirichlet}
 (\Delta_q-3)v=-1\quad\hbox{in }\Omega,\qquad
 v=c\quad\hbox{on }S,\qquad v=0\quad\hbox{on }\Sigma
\end{equation}
satisfies
\begin{equation}\label{hw-static-deficit}
 \frac23\int_\Sigma\frac{V}{H}
       -\int_\Omega V-c\kappa A
 \geq\frac32\int_\Omega
       V\left|\Hess_qv-v(\Ric_q+3q)+\frac13q\right|^2\geq0.
\end{equation}
The measures in the surface and volume integrals are those of $q$.
\end{theorem}

\begin{proof}
The static equations imply $R_q=-6$.  At $S$ they give
$\Hess_qV=0$, so $S$ is totally geodesic and $\kappa$ is a positive
constant.  If $K_S$ denotes its intrinsic Gauss curvature and $\eta$
its inward-to-$M$ unit normal, the Gauss equation gives
\[
 \Ric_q(\eta,\eta)=-3-K_S,
 \qquad (\Ric_q+3q)(\eta,\eta)=-K_S.
\]

Initially solve \eqref{hw-deficit-dirichlet} with an arbitrary
positive boundary value $c$.  The operator $-\Delta_q+3$ is
coercive with Dirichlet boundary conditions, and elliptic
regularity gives a unique smooth solution.  The maximum principle
gives $0<v\leq\max(c,1/3)$ away from $\Sigma$.
On the interior put
\[
 p=\nabla v-\frac vV\nabla V,\qquad
 B=\Hess_qv-v(\Ric_q+3q),\qquad T=B+\frac13q.
\]
Thus $\operatorname{tr}_qB=-1$ and $T$ is trace free.  Commuting
derivatives and using the static equations gives
\[
 \operatorname{div}_q B=0,\qquad
 \nabla_i p_j=B_{ij}-p_i\frac{\nabla_jV}{V},\qquad
 \operatorname{div}_q(Vp)=-V.
\]
The symmetry of $T$ therefore cancels the two terms involving
$dV$ in the following divergence:
\begin{equation}\label{hw-deficit-divergence}
 \operatorname{div}_q\bigl(VT(p,\cdot)^{\sharp}\bigr)
       =V\langle T,B\rangle_q=V|T|_q^2.
\end{equation}
The vector field on the left extends smoothly to $S$, since
$Vp=V\nabla v-v\nabla V$ does.

Write
\[
 W=\int_\Omega V,\quad
 I=W+c\kappa A,\quad
 J=\int_\Sigma\frac{V}{H},\quad
 K=\int_\Sigma VH(\partial_\nu v)^2,\quad
 E=\int_\Omega V|T|^2,
\]
where $\nu$ is the outward normal on $\Sigma$.
Green's identity gives
$I=-\int_\Sigma V\partial_\nu v$.
On $\Sigma$, the tangential Hessian is
$(\partial_\nu v)\mathrm{II}$, and hence
$T(\nu,\nu)=-2/3-H\partial_\nu v$.
On $S$, the constant boundary datum and total geodesy give
$\Hess_qv(\eta,\eta)=3c-1$, so that
$T(\eta,\eta)=c(3+K_S)-2/3$.
The outward normal of $\Omega$ along $S$ is $-\eta$.
Integrating \eqref{hw-deficit-divergence}, with these orientations,
therefore yields
\begin{equation}\label{hw-deficit-identity}
 E=\frac23 I-K+c\kappa\left(cD-\frac23 A\right),
 \qquad K+E=\frac23W+c^2\kappa D.
\end{equation}
Here Gauss--Bonnet was used to integrate $K_S$.

The weighted Cauchy--Schwarz inequality gives $I^2\leq JK$.
If $D\leq0$, the second identity in
\eqref{hw-deficit-identity} would bound $K$ above by $2W/3$
for every $c>0$, whereas $I^2/J$ tends to infinity as
$c\to\infty$.  Thus $D>0$.
Now choose $c=2A/(3D)$.  The horizon term in the first identity
vanishes, giving $E=2I/3-K$ and
\[
 J E\leq\frac23 JI-I^2
       =I\left(\frac23J-I\right).
\]
Since $I>0$, this also implies $J/I\geq3/2$ and proves
\eqref{hw-static-deficit}.
\end{proof}

\begin{lemma}[The normalized end and its boundary integral]
\label{hw-normalized-end}
In the normalized geodesic chart of
Proposition~\ref{st-end-regularity}, put $r=z^{-1}-z/4$ and
$b=dr^2/(1+r^2)+r^2\sigma$.  Then
\begin{equation}\label{hw-public-asymptotics}
 \sum_{j=0}^3|\nabla_b^j(h-b)|_b=O(r^{-3}),\qquad
 \lambda=\sqrt{1+r^2}+O_1(r^{-2}).
\end{equation}
Let $\Sigma_z$ be a small positive constant-$z$ surface, with
outward normal $\nu$ and mean curvature $H_z$.  Then
\begin{equation}\label{hw-boundary-mass-limit}
 \lim_{z\downarrow0}\int_{\Sigma_z}
       \left(\frac{2\lambda}{H_z}-\partial_\nu\lambda\right)
       =-4\pi m_{\mathrm{AH}}.
\end{equation}
\end{lemma}

\begin{proof}
Write $t=\operatorname{tr}_\sigma H_3$ and $\ell=\lambda_3$ in
\eqref{st-normalized-expansion}.  The trace constraint
\eqref{st-free-tt-data} is $t+2\ell=0$.
The exact identities
\[
 \sqrt{1+r^2}=z^{-1}+\frac z4,\qquad
 \frac{dr^2}{1+r^2}=z^{-2}dz^2,\qquad
 r^2\sigma=z^{-2}(1-z^2/4)^2\sigma
\]
and the smooth compactified expansion prove
\eqref{hw-public-asymptotics}; background orthonormal
differentiation preserves the indicated orders.
The same expansion, with $\nu=-z\partial_z$, gives
\begin{align*}
 dA_h&=z^{-2}\bigl(1-\tfrac12z^2+\tfrac12t z^3+O(z^4)\bigr)dA_\sigma,\\
 H_z&=2+z^2-\tfrac32t z^3+O(z^4),\\
 \partial_\nu\lambda&=z^{-1}-\tfrac14z-2\ell z^2+O(z^3),\\
 \frac{2\lambda}{H_z}
       &=z^{-1}-\tfrac14z+(\ell+\tfrac34t)z^2+O(z^3).
\end{align*}
In particular these spheres are strictly mean convex, and their
boundary integral tends to
$\int_{S^2}(3\ell+3t/4)dA_\sigma=-3\int_{S^2}t\,dA_\sigma/4$.
Substitution of the angular error $z^3H_3+O(z^4)$, in physical
components, into the metric-flux definition gives
\[
 p_0(h)=\frac{3}{16\pi}\int_{S^2}t\,dA_\sigma.
\]
Indeed its three leading contributions from
$V_0(\operatorname{div}_b e-d\operatorname{tr}_b e)$ and
$(\operatorname{tr}_b e)dV_0-e(\nabla_bV_0,\cdot)$ sum to
$3r^{-2}t+O(r^{-3})$ in the radial flux density.
The boundary-identity chart comparison in
Proposition~\ref{st-end-regularity} and
\eqref{st-base-equality-data} identify this flux with the original
$m_{\mathrm{AH}}$.  This proves \eqref{hw-boundary-mass-limit}.
\end{proof}

\begin{theorem}[Global identification of the equality base]
\label{hw-static-rigidity}
The completed static base is globally isometric, including its smooth
boundary and its normalized lapse, to
\begin{equation}\label{hw-global-model}
 h=\frac{dr^2}{F_M(r)}+r^2\sigma,\qquad
 \lambda=\sqrt{F_M(r)},\qquad
 F_M(r)=1+r^2-\frac{2M}{r},\qquad r\geq r_h,
\end{equation}
where the boundary is interpreted in spatial proper distance and
\begin{equation}\label{hw-actual-parameter}
 M=\frac{r_h+r_h^3}{2}=m_{\mathrm{AH}}>0.
\end{equation}
The boundary is the entire model horizon sphere, and the unique end
corresponds to $r\to\infty$.
\end{theorem}

\begin{proof}
For this proof write $V=\lambda$, $A=4\pi r_h^2$,
$\chi=\chi(S)$, and $m=m_{\mathrm{AH}}$.
Proposition~\ref{st-base-equality} supplies a smooth connected
orientable static base, complete with its compact connected
nondegenerate boundary included, and with one spherical end.
Its actual area, surface gravity, and mass satisfy
\begin{equation}\label{hw-equality-surface-gravity}
 \kappa=\frac{1+3r_h^2}{2r_h},\qquad
 A=4\pi r_h^2,\qquad m=\frac{r_h+r_h^3}{2}.
\end{equation}

\emph{Vanishing of the static deficit.}
Apply Theorem~\ref{hw-public-rigidity} on the compact domains
$\Omega_z$ between $S$ and the surfaces $\Sigma_z$ of
Lemma~\ref{hw-normalized-end}.  In particular
$D=3A+2\pi\chi>0$.  Since
\[
 3\int_{\Omega_z}V
       =\int_{\Sigma_z}\partial_\nu V-\kappa A,
\]
the left side of \eqref{hw-static-deficit} has limit
\begin{equation}\label{hw-limit-deficit}
 \frac{-4\pi m+\kappa A-2\kappa A^2/D}{3}
       =\frac{4\pi r_h^3(\chi-2)}{3(6r_h^2+\chi)}.
\end{equation}
The boundary $S$ is an orientable connected closed surface, so
$\chi\leq2$.  Its positive denominator and the nonnegative
deficit force $\chi=2$.  Thus $S$ is a sphere, and the limit in
\eqref{hw-limit-deficit} is zero.

Let $v_z$ solve \eqref{hw-deficit-dirichlet}.  Its boundary constant
$c=2A/(3D)$ is independent of $z$.  The maximum principle gives
$0<v_z\leq\max(c,1/3)$ away from the outer boundary.
Interior and boundary elliptic estimates on fixed compact subsets,
including the fixed smooth horizon, give a subsequence converging
smoothly on every such subset to a bounded function $v$.
The strong maximum principle gives $v>0$ in the interior,
and $v=c$ on $S$.  Since the right side of
\eqref{hw-static-deficit} tends to zero, its limit satisfies
\begin{equation}\label{hw-zero-tensor}
 \Hess_hv=v(\Ric_h+3h)-\frac13h
\end{equation}
throughout the interior and, by smoothness, up to $S$.

\emph{The quotient foliation.}
Put $\varphi=v/V$ and $\widehat h=V^{-2}h$ in the interior.
The conformal Hessian formula, applied to
\eqref{hw-zero-tensor} and the static equation for $V$, gives
\begin{equation}\label{hw-concircular-quotient}
 \Hess_{\widehat h}\varphi
       =\frac{\Delta_{\widehat h}\varphi}{3}\widehat h.
\end{equation}
Here $\varphi\to+\infty$ uniformly at $S$, since $v=c>0$
and $V$ has a simple zero, while $\varphi\to0$ at the unique
end because $v$ is bounded and $V\to\infty$.
Consequently $\varphi$ is proper as a map to $(0,\infty)$.
For all sufficiently large $a$, $\{\varphi\geq a\}$ is a
connected horizon collar and $\varphi=a$ is a regular level:
in horizon distance $s$, one has
$\varphi=c/(\kappa s)+O(1)$ and
$\partial_s\varphi=-c/(\kappa s^2)+O(1)$.

There are no critical points of $\varphi$.
Indeed
\[
 V\Delta_h\varphi+2\langle\nabla V,\nabla\varphi\rangle_h=-1.
\]
At a critical point, \eqref{hw-concircular-quotient} therefore
reads
$\Hess_{\widehat h}\varphi=-V\widehat h/3$.
Every critical point would be an isolated nondegenerate maximum.
Fix a high regular level in the horizon collar.  Each compact band
between it and a positive lower regular level has finitely many
critical points.  On decreasing the level through a nondegenerate
maximum the superlevel set acquires a disjoint ball; between
critical levels the normalized gradient flow preserves its
components.  No such components can merge without a different
kind of critical point.  A component created at a maximum could
thus never join the horizon component at a positive lower level.
This contradicts a path from that maximum to the horizon collar,
whose positive continuous function $\varphi$ has a positive
minimum.  This proves the assertion.
Normalized gradient flow on compact bands now gives the global
product by connected spherical levels of $\varphi$.

\emph{Warped splitting on a collar and on the whole base.}
Fix a large $a$ and write $\Sigma_a=\{\varphi=a\}$ and
$\Omega_a=\{\varphi\geq a\}\cup S$.
The function $w=v-aV$ equals $c$ on $S$, is zero on $\Sigma_a$,
is positive inside, and satisfies \eqref{hw-zero-tensor} with
$v$ replaced by $w$.  On $\Sigma_a$, its tangential equation is
\[
 (\partial_\nu w)\mathrm{II}=-\frac13h|_{T\Sigma_a},
 \qquad H\partial_\nu w=-\frac23.
\]
Since $\partial_\nu w<0$, the surface is strictly mean convex.
Green's identity consequently gives the exact finite-domain
equality
\[
 \frac23\int_{\Sigma_a}\frac{V}{H}
       =-\int_{\Sigma_a}V\partial_\nu w
       =\int_{\Omega_a}V+c\kappa A.
\]
The hypotheses of
\cite[Theorem~1.1]{borghinifogagnolopinamonti2024} all hold:
the substatic tensor is zero, $S$ and $\Sigma_a$ are connected,
and $\Hess_hV/V=\Ric_h+3h$ is smooth up to $S$.
Its horizon constant in Equation~(1.4) is exactly
$2A/(3D)=c$.  That theorem makes this entire collar a warped
product with radial $V$.  The solution $w$ is radial there as well:
its average over a warped-product fiber solves the same
$\Delta_h-3$ Dirichlet problem with the same constant boundary
data, and uniqueness identifies that average with $w$.
Thus this collar splitting agrees with the $\varphi$ foliation.

Equation~\eqref{hw-concircular-quotient}, and the connectedness
of the levels, give a global optical warped product
\[
 \widehat h=d\rho^2+\psi(\rho)^2\gamma.
\]
To see this directly, tangential differentiation shows that
$|\nabla\varphi|_{\widehat h}^2$ is a function of $\varphi$;
its derivative determines the pure-trace Hessian coefficient.
The unit normal flow therefore rescales the level metric by a
factor depending only on its normal coordinate $\rho$.
For $W=1/V$, the static equations in this conformal metric become
\[
 \Hess_{\widehat h}W=-\frac W2\Ric_{\widehat h}.
\]
Its radial component is
$\partial_\rho^2W=(\psi''/\psi)W$.
Both initial data $W$ and $\partial_\rho W$ are constant on a
fiber in the collar just obtained.  Uniqueness of this scalar
ordinary differential equation makes $W$, and hence $V$, radial
on the whole product.  Replacing $\rho$ by physical distance
therefore gives $h=dt^2+R(t)^2\gamma$ globally.

\emph{Identification and normalization.}
The scalar equation $R_h=-6$ forces the Gauss curvature of
$\gamma$ to be constant.  Since its underlying surface is $S^2$,
this curvature is positive.  Rescale the fiber to the unit round
metric and write the metric as $dt^2+r(t)^2\sigma$.
At the horizon $r(0)=r_h$ and $r'(0)=0$.
The scalar equation and its first integral are
\[
 2rr''+(r')^2=1+3r^2,\qquad
 \frac r2\bigl(1+r^2-(r')^2\bigr)=M.
\]
Thus $M=(r_h+r_h^3)/2>0$ and
$r''(0)=(1+3r_h^2)/(2r_h)>0$.
It follows that $r'>0$ initially and then everywhere outside
the horizon: the first integral gives $(r')^2=F_M(r)$, whose
strictly increasing function $F_M$ has the unique positive zero
$r_h$.
The tangential static equation gives
$V'r'=Vr''$, so $V=C r'$ for one positive constant $C$.
At the horizon, \eqref{hw-equality-surface-gravity} gives
$\kappa=V'(0)=C r''(0)$, whence $C=1$.
The unique end has $V\to\infty$; hence $r\to\infty$, and the
global product covers the full model exterior.
These identities prove \eqref{hw-global-model} with the same
unrescaled lapse and with the original mass
\eqref{hw-actual-parameter}.

The product isometry includes the full horizon.  Its extension to
the completed boundary is smooth in proper distance: on a
positive lapse level it is smooth and intertwines normal geodesic
transport to the two smooth nondegenerate horizons.
Finally, the argument divides neither by $3r_h^2-1$ nor by
$\kappa^2-3$.  It applies at $r_h=1/\sqrt3$, where
$\kappa=\sqrt3$, and on both sides of that value.  Equivalently,
both roots of $3r_h^2-2\kappa r_h+1=0$ are covered.
\end{proof}
